% Generated draft: requires mathematical and source review.
\documentclass[11pt]{article}
\usepackage[margin=25mm]{geometry}
\usepackage{fontspec}
\setmainfont{Latin Modern Roman}
\usepackage{amsmath,amssymb,mathtools,unicode-math}
\setmathfont{Latin Modern Math}
\usepackage{enumitem,needspace,xurl,hyperref}
\setlength{\parindent}{0pt}
\setlength{\parskip}{5pt}
\setlength{\emergencystretch}{3em}
\newcommand{\N}{\mathbb{N}}
\newcommand{\Z}{\mathbb{Z}}
\newcommand{\Q}{\mathbb{Q}}
\newcommand{\R}{\mathbb{R}}
\newcommand{\C}{\mathbb{C}}
\newcommand{\F}{\mathbb{F}}
\newcommand{\A}{\mathbb{A}}
\newcommand{\PP}{\mathbb{P}}
\newcommand{\E}{\mathbb{E}}
\newcommand{\eps}{\varepsilon}
\newcommand{\dd}{\,\mathrm{d}}
\begin{document}
\section*{UnsolvedMath ID: 10400152}
\textbf{Source problem code:} catalogue entry
\subsection{Definitions and mathematical statement}
An HQFT is a homotopy quantum field theory: in a specified dimension and with a fixed target space \(X\), it assigns algebraic state objects to manifolds equipped with maps to \(X\), and linear maps to the corresponding cobordisms, compatibly with gluing, disjoint union and the prescribed homotopy equivalence of target maps. The source attributes two tasks to V. Turaev:
\begin{enumerate}
\item Extend the HQFT framework to spin and \(\operatorname{spin}^c\) settings. One must specify bordisms carrying the relevant tangential structure together with maps to \(X\), their equivalences, and field-theory axioms compatible with those structures. A spin structure is a lift of the oriented frame bundle through \(\operatorname{Spin}(d)\to SO(d)\); a \(\operatorname{spin}^c\) structure uses \(\operatorname{Spin}^c(d)=(\operatorname{Spin}(d)\times U(1))/\{(1,1),(-1,-1)\}\). Boundary structures and their compatibility under gluing must be included.
\item Find the algebraic structures underlying spin and \(\operatorname{spin}^c\) HQFTs in dimension \(1+1\): describe the data and relations that encode the state assignments to structured one-manifolds and the operations of structured two-dimensional cobordisms with maps to \(X\), with a reconstruction/equivalence statement in the chosen HQFT category.
\end{enumerate}
For example, after the missing conventions are fixed, a functorial formulation would involve symmetric monoidal functors from the appropriate structured bordism category to a category of modules or vector spaces. This describes the kind of framework sought, not a selected ground ring, a claimed classification or an assertion that every oriented HQFT admits a refinement. The first task is not explicitly restricted to \(1+1\) dimensions; that restriction belongs to the second task. The source does not identify the homotopy target, coefficient category, finiteness conditions or algebraic equivalence relation.
\subsection{Short English statement}
Turaev asks to formulate HQFTs with spin and \(\operatorname{spin}^c\) structures and to identify the algebraic data that describe them in dimension \(1+1\). The homotopy target and coefficient conventions remain unspecified.
\subsection{Sources}
\begin{itemize}
\item Local UnsolvedMath registry, permanent ID \(10400152\), exact source code \texttt{catalogue entry}.
\item V. Turaev, attribution in Ohtsuki, \emph{Problems on invariants of knots and 3-manifolds} (2002), Problem 8.6, PDF page 124: \url{https://msp.org/gtm/2002/04/gtm-2002-04-024s.pdf}. Registry access 2026-07-29, extraction pdf-font-structure.
\end{itemize}
\end{document}
